% answers_short.tex —— 1994 年·全国卷·理科 · 简短答案速查
% 由 tools/make_short_answers.py 自 answers.tex 的 \gksolhead/\gkanswer 提取，
% 并与 1994/1994_数学真题_全国理_答案要点.md 逐题比对后生成，勿手改（改源 MD 后重跑）。

\begin{gknotice}
\textbf{使用说明}：本节只给最终答案便于快速核对，完整推导见解析版。答案为编写者独立推导，\textbf{非官方标准答案}。
\end{gknotice}

\gkpart{选择题}
\begin{center}\begin{tabular}{cccccccccccccccc}
\toprule
题号 & 1 & 2 & 3 & 4 & 5 & 6 & 7 & 8 & 9 & 10 & 11 & 12 & 13 & 14 & 15 \\
\midrule
答案 & C & D & D & A & B & D & B & A & A & C & C & B & D & B & C \\
\bottomrule
\end{tabular}\end{center}

\gkpart{填空题}
\begin{itemize}[leftmargin=2.6em]
  \item[\textbf{16.}] \(-189\)
  \item[\textbf{17.}] 准线 \(x=3\)；圆 \((x-2)^2+y^2=1\)
  \item[\textbf{18.}] \(-\dfrac34\)
  \item[\textbf{19.}] \(\dfrac{2\sqrt2}{3}\pi\)
  \item[\textbf{20.}] \(\dfrac{a_1+a_2+\cdots+a_n}{n}\)
\end{itemize}

\gkpart{解答题}
\begin{itemize}[leftmargin=2.6em]
  \item[\textbf{21.}] (1) \(\omega=\sqrt2\left[\cos\left(-\dfrac{3\pi}{4}\right)+i\sin\left(-\dfrac{3\pi}{4}\right)\right]\)；(2) \(a=-1,\ b=2\)
  \item[\textbf{22.}] 见解析（\(s,t\) 代换＋和差化积）
  \item[\textbf{23.}] (1) 见解析（\(DM\parallel AB_1\)，中位线）；(2) \(45^\circ\)
  \item[\textbf{24.}] \(l:\ y=\dfrac{1+\sqrt5}{2}x\)；\(C:\ y^2=\dfrac{4\sqrt5}{5}x\)
  \item[\textbf{25.}] (1) \(a_1=2,\ a_2=6,\ a_3=10\)；(2) \(a_n=4n-2\)；(3) \(1\)
\end{itemize}
